\section{The finite-interval action inequality}\label{sec:action}

Fix
\[
 0<\sigma<1,\qquad p=1+\sigma,\qquad q=1+\frac1\sigma.
\]
All matrices below are finite dimensional. Write $\Sym_N$ for the real
symmetric $N\times N$ matrices, with inner product $\tr(AB)$ and norm
$\|A\|_{\HS}=(\tr(A^2))^{1/2}$. The notation $A\le B$ means that $B-A$
is positive semidefinite. Vectors are columns, and $|a|$ is their Euclidean norm.

Let $I=(x_-,x_+)$ be a bounded interval. For a positive diagonal matrix
$K=\diag(k_1,\ldots,k_N)$ and $a\in H^1_0(I;\R^N)$, define
\begin{equation}\label{eq:action}
 \mathcal E_K(a)=\int_I\bigl(|a'|^2+a^{\transpose}K^2a-|a|^{2q}\bigr)\,dx.
\end{equation}
In one dimension $H^1_0(I)$ embeds into the continuous functions, so every
term in this integral is finite.

\begin{theorem}[Action inequality]\label{thm:action}
Let $u=(u_1,\ldots,u_N)^{\transpose}\in H^1_0(I;\R^N)$ satisfy
$\int_I uu^{\transpose}\,dx=I_N$, and let $k_i>0$. Then
\begin{equation}\label{eq:action-bound}
 \sup_{C\text{ lower triangular}}\mathcal E_K(Cu)
 \ge \sum_{i=1}^N\int_{-k_i}^{k_i}(k_i^2-s^2)^\sigma\,ds.
\end{equation}
The supremum is over all real lower triangular $N\times N$ matrices.
\end{theorem}

The restriction to lower triangular matrices is what makes this inequality
useful spectrally. If the $u_i$ are eigenfunctions ordered from the lowest
eigenvalue upward, the $i$th component of $Cu$ stays in the span of the first
$i$ eigenfunctions. Section~\ref{sec:spectral} will use this observation to
bound the left side from above by an integral of the potential.

Here is the construction behind the lower bound. For each $\delta>0$,
Section~\ref{sec:field} builds a locally Lipschitz map
$M:\Sym_N\to\Sym_N$ and a $C^1$ function $J:\Sym_N\to\R$ such that
\begin{equation}\label{eq:framework-gradient}
 dJ(B)[H]=-\tr(M(B)H).
\end{equation}
They have two decisive properties. First, Section~\ref{sec:rankone} proves
\begin{equation}\label{eq:framework-rankone}
 M(B)=aa^{\transpose}\quad\Longrightarrow\quad
 a^{\transpose}(K^2-B^2)a\ge |a|^{2q}.
\end{equation}
Second, Section~\ref{sec:path} finds, for every $\varepsilon>0$, a lower
triangular $C$ and a path $B$ satisfying
\begin{equation}\label{eq:framework-path}
 \varepsilon B'=M(B)-(Cu)(Cu)^{\transpose},\qquad
 B(x_-)=K,\quad B(x_+)=-K.
\end{equation}
The small parameter $\varepsilon$ penalizes failure of the rank-one constraint.
For $a=Cu$, differentiating $J(B)+a^{\transpose}Ba$ along this path produces
the negative term $-\varepsilon^{-1}\|M(B)-aa^{\transpose}\|_{\HS}^2$.
Uniform bounds on the path and on $C$ then give
\[
 \sup_C\mathcal E_K(Cu)\ge J(-K)-J(K).
\]
Finally, the boundary difference tends to the right side of
\eqref{eq:action-bound} as $\delta\downarrow0$. We keep $\delta$ fixed until
after the limit $\varepsilon\downarrow0$.
